
% 2. Analiza Rozmieszczenia (Placement)
% • Badanie wpływu fizycznej lokalizacji oscylatorów na krzemie na stopień ich
% niezależności i czystą entropię.
% • Przygotowanie konfiguracji z różną liczbą oscylatorów pierścieniowych (RO): 8, 16,
% 64 oraz 128 jednostek. (dobrać wartości tak aby uchwycić jakąś różnicę)
% • Realizacja dwóch skrajnych wariantów rozmieszczenia: Localized (bloki RO
% upakowane ciasno w jednym obszarze CLB) oraz Distributed (rozmieszczenie
% rozproszone po całej strukturze FPGA przy użyciu ograniczeń
% projektowych/constraints).
% • Pomiary bazowe: zużycie zasobów logicznych (liczba tablic LUT), średni cykl
% oscylacji poszczególnych RO, Min-entropia ($H_{min}$) oraz entropia Shannona
% (pakiet ent).
% • TABELA 1: Wpływ fizycznego rozmieszczenia na źródło entropii.
% • Zawartość: Liczba RO | Typ rozmieszczenia (Loc/Dist) | Liczba LUT | Wynik
% $H_{min}$ | Entropia wg ent.
% • Cel: Wykazanie występowania korelacji przestrzennej i jej wpływu na jakość
% generowanego szumu.

\chapter{Analiza rozmieszczenia}